\subsection{Laplace Equation}
Describes the pressure difference across a curved interface.
\begin{equation}
    \Delta P = \gamma \left( \frac{1}{R_1} + \frac{1}{R_2} \right)
\end{equation}
\textbf{Special Cases:}
\begin{itemize}
    \item \textbf{Sphere:} $\Delta P = \frac{2\gamma}{R}$
    \item \textbf{Cylinder:} $\Delta P = \frac{\gamma}{R}$
    \item \textbf{Planar:} $\Delta P = 0$
\end{itemize}

\subsection{Meniscus and Capillarity}
\begin{itemize}
    \item \textbf{Concave Meniscus:} Adhesion $>$ Cohesion ($r < 0$).
    \item \textbf{Convex Meniscus:} Cohesion $>$ Adhesion ($r > 0$).
\end{itemize}
\textbf{Ostwald Ripening:} Process where large particles/droplets grow at the expense of smaller ones to minimize total surface energy.

\subsection{Kelvin Equation}
Describes the change in vapor pressure due to a curved liquid-vapor interface (droplet or bubble).
\begin{equation}
    \ln \left( \frac{P_r}{P^0} \right) = \frac{2 \gamma V_m}{RT r}
\end{equation}
Where $P_r$ is the vapor pressure over the curved surface, and $P^0$ is the standard vapor pressure.
For \textbf{Capillary Condensation} in a pore with radius $r$:
\begin{equation}
    RT \ln \left( \frac{P}{P_0} \right) = \gamma V_m \left( \frac{1}{R_1} + \frac{1}{R_2} \right)
\end{equation}

\subsection{Nucleation}
\textbf{Critical Radius ($r_c$):}
\begin{equation}
    r_c = \frac{2 \gamma V_m}{RT \ln S}
\end{equation}
\textbf{Critical Free Energy ($\Delta G_c$):}
\begin{equation}
    \Delta G_c = \frac{16 \pi \gamma^3 V_m^2}{3 (RT \ln S)^2}
\end{equation}

\subsection{Marangoni Effect}
Tangential (shear) forces created by a gradient in surface tension along an interface:
\begin{equation}
    f_{Marangoni} = \nabla \gamma
\end{equation}
Causes liquid to flow from regions of low surface tension to high surface tension.
\vspace{0.2cm}